\documentclass[../../main.tex]{subfiles}
\begin{document}
\section{接触结构}\label{section:ch22-contact-structures}

正如我们所见，辛流形必须是偶维的；但有一种与之密切相关的结构，称为\emph{接触结构}，可以在奇维流形上定义。它在几何与分析中也有重要应用。本节中，我们介绍接触几何的主要内容。

设 $M$ 是维数为 $2n+1$ 的光滑流形。$M$ 上的一个\emph{接触形式}是一个处处非零的光滑 $1$-形式 $\tau$，满足：对每个 $p\in M$，$\mathrm{d}\tau_p$ 在子空间 $\ker\tau_p\subset T_pM$ 上的限制是非退化的，也就是说是一个辛张量。$M$ 上的一个\emph{接触结构}是一个秩为 $2n$ 的光滑分布 $H\subset TM$，满足 $H$ 的任意光滑局部定义形式都是接触形式。一个\emph{接触流形}是 $M$ 连同 $M$ 上的一个接触结构。若 $(M,H)$ 是接触流形，$H$ 的任意（局部或整体）定义形式称为 $H$ 的一个接触形式。Jean Martinet 于 1971 年证明 [Mar71]：每个可定向紧光滑 $3$ 流形都容许接触结构；但哪些更高维流形容许接触结构这一问题仍未解决。

\begin{proposition}\label{proposition:ch22-contact-form-criterion}
$(2n+1)$ 流形 $M$ 上的光滑 $1$-形式 $\tau$ 是接触形式，当且仅当 $\tau\wedge(\mathrm{d}\tau)^n$ 在 $M$ 上处处非零；其中 $(\mathrm{d}\tau)^n$ 表示 $n$ 重外积 $\mathrm{d}\tau\wedge\cdots\wedge\mathrm{d}\tau$。
\end{proposition}

\begin{exercise}\label{exercise:ch22-25}
证明前一个命题。
\end{exercise}

\begin{exercise}\label{exercise:ch22-26}
设 $H$ 是光滑流形 $M$ 上的接触结构。证明：若 $\tau_1$ 和 $\tau_2$ 是 $H$ 的任意两个局部接触形式，则在它们的公共定义域上存在光滑处处非零函数 $f$ 使 $\tau_2=f\,\tau_1$。
\end{exercise}

由习题 19-2 的结果可知，余维 $1$ 的分布 $H\subset TM$ 是可积的，当且仅当任意局部定义形式 $\tau$ 满足 $\tau\wedge\mathrm{d}\tau\equiv0$。相比之下，若 $H$ 是接触结构，则不仅 $\tau\wedge\mathrm{d}\tau$ 在 $M$ 上处处非零，而且再与 $\mathrm{d}\tau$ 作 $n-1$ 次外积后仍非零。因此，接触结构在某种意义上是一个“极大不可积分布”。

\begin{example}[接触形式]\label{example:ch22-contact-forms}
\begin{enumerate}[(a)]
\item 在带坐标 $(x^1,\dots,x^n,y^1,\dots,y^n,z)$ 的 $\mathbb{R}^{2n+1}$ 上，定义 $1$-形式 $\tau$ 为
\begin{align}
\tau=\mathrm{d}z-\sum_{i=1}^n y^i\,\mathrm{d}x^i,
\end{align}
并令 $H\subset T\mathbb{R}^{2n+1}$ 为被 $\tau$ 零化的秩 $2n$ 分布。于是 $\mathrm{d}\tau=\sum_{i=1}^n\mathrm{d}x^i\wedge\mathrm{d}y^i$。若我们定义向量场族 $\{X_i,Y_i:i=1,\dots,n\}$ 为
\begin{align*}
X_i=\frac{\partial}{\partial x^i}+y^i\frac{\partial}{\partial z},\qquad
Y_i=\frac{\partial}{\partial y^i},
\end{align*}
则 $\{X_i,Y_i\}$ 是 $H$ 的一个光滑标架，并且直接验证可知它满足 $\mathrm{d}\tau(X_i,X_j)=\mathrm{d}\tau(Y_i,Y_j)=0$ 及 $\mathrm{d}\tau(X_i,Y_j)=\delta_{ij}$。正如在\cref{example:ch22-standard-symplectic-tensor}中一样，由此可知 $\mathrm{d}\tau|_H$ 非退化，故 $\tau$ 是接触形式。下面的\cref{theorem:ch22-contact-darboux}表明：每个接触形式都可局部地通过坐标变换化为这一形式。
\item 设 $M$ 是光滑 $n$ 流形，并在 $(2n+1)$ 流形 $\mathbb{R}\times T^*M$ 上定义光滑 $1$-形式 $\tau=\mathrm{d}z-\sigma$，其中 $z$ 是 $\mathbb R$ 上的标准坐标，$\sigma$ 是 $T^*M$ 上的重言 $1$-形式。用 $T^*M$ 的自然坐标 $(x^i,\xi_i)$，该形式的坐标表示为
\begin{align*}
\tau=\mathrm{d}z-\sum_{i=1}^n \xi_i\,\mathrm{d}x^i,
\end{align*}
故由与 (a) 相同的论证它是接触形式。
\item 在带坐标 $(x^1,\dots,x^{n+1},y^1,\dots,y^{n+1})$ 的 $\mathbb{R}^{2n+2}$ 上，考虑 $1$-形式
\begin{align*}
\rho=\sum_{i=1}^{n+1}(x^i\,\mathrm{d}y^i-y^i\,\mathrm{d}x^i).
\end{align*}
$S^{2n+1}$ 上的标准接触形式是光滑 $1$-形式 $\tau=\iota^*\rho$，其中 $\iota:S^{2n+1}\hookrightarrow\mathbb{R}^{2n+2}$ 是包含。为看清 $\tau$ 确实是接触形式，首先注意 $\mathrm{d}\rho=2\sum_{i=1}^{n+1}\mathrm{d}x^i\wedge\mathrm{d}y^i$ 是 $\mathbb{R}^{2n+2}$ 上的辛形式。考虑 $\mathbb{R}^{2n+1}\setminus\{0\}$ 上的下列向量场
\begin{align*}
N=x^j\frac{\partial}{\partial x^j}+y^j\frac{\partial}{\partial y^j},\qquad
T=x^j\frac{\partial}{\partial y^j}-y^j\frac{\partial}{\partial x^j}.
\end{align*}
一个计算表明 $N$（关于 Euclid 度量）垂直于 $S^{2n+1}$，而 $T$ 切于它。设 $S\subset T(\mathbb{R}^{2n+2}\setminus\{0\})$ 表示由 $N$ 和 $T$ 张成的子丛，并设 $S^\omega$ 表示它关于 $\mathrm{d}\rho$ 的辛补。对每个 $p\in S^{2n+1}$，$S^\omega_p$ 是满足 $\mathrm{d}\rho_p(N_p,X_p)=\mathrm{d}\rho_p(T_p,X_p)=0$ 的向量 $X\in T_p\mathbb{R}^{2n+2}$ 的集合。我们计算
\begin{align*}
N\lrcorner\,\mathrm{d}\rho=2\sum_{i=1}^{n+1}(x^i\,\mathrm{d}y^i-y^i\,\mathrm{d}x^i)=2\rho,
\end{align*}
\begin{align*}
T\lrcorner\,\mathrm{d}\rho=2\sum_{i=1}^{n+1}(x^i\,\mathrm{d}x^i+y^i\,\mathrm{d}y^i)=\mathrm{d}(|x|^2+|y|^2).
\end{align*}
由此可知 $S^\omega_p$ 是 $\rho_p$ 与 $\mathrm{d}(|x|^2+|y|^2)_p$ 的公共核，即 $\ker\rho_p\cap T_pS^{2n+1}=\ker\tau_p$。由于在 $\mathbb{R}^{2n+1}\setminus\{0\}$ 上 $\mathrm{d}\rho(N,T)=|x|^2+|y|^2\neq0$，$S_p$ 是 $T_p\mathbb{R}^{2n+2}$ 的辛子空间，从而由\cref{proposition:ch22-subspace-types}(a)，$\ker\tau_p=S^\omega_p$ 也是辛子空间。由于 $\mathrm{d}\tau_p$ 在 $\ker\tau_p$ 上的限制与 $\mathrm{d}\rho_p$ 的限制相同，它是非退化的，故 $\tau$ 是接触形式。
\end{enumerate}
\end{example}

\begin{theorem}[Reeb 场]\label{theorem:ch22-reeb}
设 $(M,H)$ 是接触流形，并设 $\tau$ 是 $H$ 的一个接触形式。存在唯一向量场 $T\in\mathfrak X(M)$，称为 $\tau$ 的 Reeb 场，满足下列两个条件：
\begin{align}
T\lrcorner\,\mathrm{d}\tau=0,\qquad \tau(T)=1.
\end{align}
\end{theorem}

\begin{proof}
定义光滑丛同态 $\Phi:TM\to T^*M$ 为 $\Phi(X)=X\lrcorner\,\mathrm{d}\tau$，并对每个 $p\in M$，令 $\Phi_p$ 表示线性映射 $\Phi|_{T_pM}:T_pM\to T^*_pM$。$\mathrm{d}\tau_p$ 在 $H_p$ 上限制为非退化 $2$-张量这一事实蕴含 $\Phi_p|_{H_p}$ 是单射，故 $\Phi_p$ 的秩至少为 $2n$（其中 $2n+1$ 是 $M$ 的维数）。另一方面，$\Phi_p$ 不可能有秩 $2n+1$，因为那样 $\mathrm{d}\tau_p$ 就会是非退化的，而这不可能，因为 $T_pM$ 是奇维的。因此 $\Phi_p$ 的秩恰为 $2n$，故 $\dim\ker\Phi_p=1$。由于 $\ker\Phi_p$ 不包含在 $H_p=\ker\tau_p$ 中，存在唯一向量 $T_p\in\ker\Phi_p$ 满足 $\tau_p(T_p)=1$。这表明存在唯一的粗向量场 $T$ 满足(22.20)。

为看清 $T$ 是光滑的，注意由定理 10.34，$\ker\Phi$ 是 $TM$ 的光滑秩 $1$ 子丛。给定 $p\in M$，令 $X$ 为 $\ker\Phi$ 在 $p$ 的某个邻域上任意光滑处处非零截面。由于 $\tau(X)\neq0$，我们可局部写 Reeb 场为 $T=\tau(X)^{-1}X$，它也是光滑的。
\end{proof}

\begin{exercise}\label{exercise:ch22-29}
证明\cref{example:ch22-contact-forms}中描述的三个接触形式的 Reeb 场如下：
\begin{enumerate}[(a)]
\item $T=\partial/\partial z$；
\item $T=\partial/\partial t$；
\item $T=\big(x^j\,\partial/\partial y^j-y^j\,\partial/\partial x^j\big)|_{S^{2n+1}}$。
\end{enumerate}
\end{exercise}

\begin{exercise}\label{exercise:ch22-30}
设 $\tau$ 是接触形式，$T$ 是其 Reeb 场。证明 $\mathcal L_T\tau=0$。
\end{exercise}

我们为辛流形描述的许多构造，在接触几何中都有类似物。我们从 Darboux 定理的一个类似物开始。

\begin{theorem}[接触 Darboux 定理]\label{theorem:ch22-contact-darboux}
设 $\tau$ 是 $(2n+1)$ 维流形 $M$ 上的接触形式。对每个 $p\in M$，存在以 $p$ 为中心的光滑坐标 $(x^1,\dots,x^n,y^1,\dots,y^n,z)$，使 $\tau$ 具有式(22.19)的形式。
\end{theorem}

\begin{proof}
任取 $p\in M$。设 $(U,(u^i))$ 是以 $p$ 为中心的光滑坐标立方体，其中 $\tau$ 的 Reeb 场具有形式 $T=\partial/\partial u^1$；并令 $Y\subset U$ 为由 $u^1=0$ 定义的切片。由于 $T$ 处处不切于 $Y$，$\mathrm{d}\tau$ 在 $Y$ 上的拉回是辛形式。必要时收缩 $U$ 和 $Y$，我们可以找到以 $p$ 为中心的 $Y$ 的 Darboux 坐标 $(x^1,\dots,x^n,y^1,\dots,y^n)$，并通过要求它们与 $u^1$ 无关（等价地，在 $T$ 的积分曲线上为常数）把它们延拓到 $U$ 上。令 $\alpha$ 为 $U$ 上的 $1$-形式 $\sum_i y^i\,\mathrm{d}x^i$，于是 $\mathrm{d}\tau$ 与 $-\mathrm{d}\alpha$ 在 $Y$ 上的拉回一致。由于 $T\lrcorner\,\mathrm{d}\tau=T\lrcorner\,\mathrm{d}\alpha=0$，可知在 $Y$ 的点处 $\mathrm{d}\tau+\mathrm{d}\alpha=0$。然后 $\mathcal L_T\tau=\mathcal L_T\alpha=0$ 蕴含 $\mathrm{d}(\tau+\alpha)$ 在 $T$ 的流下不变，故事实上 $\mathrm{d}(\tau+\alpha)=0$ 在整个 $U$ 上成立。由 Poincar\'e 引理，存在 $U$ 上的光滑函数 $z$ 使 $\mathrm{d}z=\tau+\alpha$；通过减去一个常数，可安排 $z(p)=0$。由于 $\mathrm{d}z_p(T_p)=\tau_p(T_p)=1$，可知 $\{\mathrm{d}x^i|_p,\mathrm{d}y^i|_p,\mathrm{d}z|_p\}$ 线性无关，故习题 11-6 表明在 $p$ 的某个邻域上 $(x^1,\dots,x^n,y^1,\dots,y^n,z)$ 就是我们所求的坐标。
\end{proof}

下一个命题描述 Hamilton 向量场的一个接触类似物。

\begin{proposition}\label{proposition:ch22-contact-hamiltonian}
设 $(M,H)$ 是接触流形，$\tau$ 是 $H$ 的接触形式。对任意函数 $f\in C^\infty(M)$，存在唯一向量场 $X_f$，称为 $f$ 的\emph{接触 Hamilton 向量场}，满足 $\tau(X_f)=f$ 且 $(X_f\lrcorner\,\mathrm{d}\tau)|_H=-\mathrm{d}f|_H$。
\end{proposition}

\begin{proof}
设 $f\in C^\infty(M)$。由于 $\mathrm{d}\tau$ 在 $H$ 上的限制非退化，存在唯一光滑向量场 $B\in\Gamma(H)$ 使 $B\lrcorner\,\mathrm{d}\tau|_H=\mathrm{d}f|_H$。若令 $X_f=fT-B$（其中 $T$ 是 $\tau$ 的 Reeb 场），则容易验证所需条件满足。
\end{proof}

设 $(M,H)$ 是接触流形。光滑向量场 $X\in\mathfrak X(M)$ 如果其流 $\theta$ 保持接触结构（即对 $\theta$ 定义域中所有 $(t,p)$ 都有 $\mathrm{d}(\theta_t)_p(H_p)=H_{\theta_t(p)}$），就称为\emph{接触向量场}。

\begin{theorem}[接触向量场的刻画]\label{theorem:ch22-contact-vector-field}
若 $(M,H)$ 是接触流形，$\tau$ 是 $H$ 的接触形式，则 $M$ 上一个光滑向量场是接触向量场，当且仅当它是接触 Hamilton 向量场。
\end{theorem}

\begin{proof}
习题 22-21。
\end{proof}

若 $(M,H)$ 是接触流形，光滑子流形 $S\subset M$ 如果 $TS\subset H$（等价地，对任意接触形式 $\tau$ 都有 $\iota^*\tau=0$，其中 $\iota:S\hookrightarrow M$ 是包含），就称为\emph{迷向的}。若 $S\subset M$ 是迷向的，则 $\iota^*\mathrm{d}\tau=\mathrm{d}(\iota^*\tau)=0$。这蕴含对每个 $p\in S$，切空间 $T_pS$ 是辛向量空间 $H_p$ 的迷向子空间，因而其维数不能大于 $n$（其中 $2n+1$ 是 $M$ 的维数）。具有最大可能维数 $n$ 的迷向子流形称为\emph{Legendrian 子流形}。

下一个定理是 Hamilton 流出定理的接触类似物，其证明方式大致相同。它是构造完全非线性 PDE 解的主要工具（见下面的\cref{theorem:ch22-general-cauchy}）。

\begin{theorem}[接触流出定理]\label{theorem:ch22-contact-flowout}
设 $(M,H)$ 是接触流形，$F\in C^\infty(M)$，$\Sigma$ 是 $M$ 的嵌入迷向子流形且包含在 $F$ 的零集中，并设接触 Hamilton 向量场 $X_F$ 处处不切于 $\Sigma$。若 $\Gamma$ 是从 $\Sigma$ 沿 $X_F$ 的流出，则 $\Gamma$ 也是迷向的，并包含在 $F$ 的零集中。
\end{theorem}

\begin{proof}
习题 22-23。
\end{proof}
\end{document}
